\chapter{Results and Discussion}
\label{cha:results}

This chapter presents the experimental results for the signal denoising
stage~(ATP-1) and the data compression stage~(ATP-3). Section~\ref{sec:nist-denoise}
reports NIST denoising results. Section~\ref{sec:nist-compress} reports NIST
compression results. Section~\ref{sec:podfam-results} presents results on the
PODFAM industrial dataset. Section~\ref{sec:discussion} discusses the findings
in relation to the two research questions.

\section{Denoising Results on NIST Dataset (ATP-1)}
\label{sec:nist-denoise}

\subsection{Quantitative Results}

Table~\ref{tab:denoise-results} reports the noise standard deviation and
percentage noise reduction for all ten denoising methods on NIST Layer~01.
The raw signal baseline has a noise standard deviation of
500.2\,\textdegree C, reflecting the significant electronic noise and infrared
reflections characteristic of the laser powder bed fusion environment.

\begin{table}[H]
\centering
\caption{Denoising results on NIST Layer~01~(2,065 frames). Raw baseline noise
         standard deviation = 500.2\,\textdegree C. Methods sorted by noise
         reduction within each category.}
\label{tab:denoise-results}
\begin{tabular}{llccc}
\toprule
\textbf{Method} & \textbf{Type} &
\textbf{Noise Std Dev (\textdegree C)} & \textbf{Reduction (\%)} & \textbf{Rank} \\
\midrule
Raw Baseline     & Baseline  & 500.2 &  0.0 & --- \\
\midrule
Kalman Filter    & Classical &  52.5 & 89.5 & 1 \\
Gaussian Filter  & Classical & 210.1 & 58.0 & 2 \\
Moving Average   & Classical & 280.5 & 43.9 & 3 \\
Savitzky--Golay  & Classical & 358.4 & 28.3 & 4 \\
Median Filter    & Classical & 363.3 & 27.4 & 5 \\
\midrule
LSTM Denoiser    & ML        & 341.3 & 31.8 & 1 \\
CNN Denoiser     & ML        & 350.3 & 30.0 & 2 \\
BiLSTM Denoiser  & ML        & 355.0 & 29.0 & 3 \\
Autoencoder Den. & ML        & 363.0 & 27.5 & 4 \\
\bottomrule
\end{tabular}
\end{table}

The Kalman filter achieves the best result by a large margin, reducing noise
by~89.5\,\% to a standard deviation of just 52.5\,\textdegree C. The Gaussian
filter is the second-best classical method~(58.0\,\%) and the moving average
third~(43.9\,\%). Among the ML methods, the LSTM performs best~(31.8\,\%),
followed by the CNN~(30.0\,\%), the BiLSTM~(29.0\,\%), and the
Autoencoder~(27.5\,\%).

Figure~\ref{fig:denoise-compare} shows the noise standard deviation bar chart
for all methods.

\begin{figure}[H]
\centering
\includegraphics[width=\textwidth]{figures/full_comparison_denoise}
\caption{Noise standard deviation for all ten denoising methods on NIST Layer~01.
         Classical methods are shown in blue/teal, ML methods in orange/red.
         The Kalman filter~(52.5\,\textdegree C) achieves noise levels far
         below all other methods.}
\label{fig:denoise-compare}
\end{figure}

\subsection{Discussion}

The Kalman filter's superiority is explained by the nature of the noise. The
NIST signal is dominated by Gaussian electronic noise on a slowly varying
temperature trend, which is exactly the signal type the Kalman filter was
designed for~\cite{kalman_wavelet_denoising}. The ratio $Q/R = 10^{-4}$ means
the filter trusts its model prediction strongly and suppresses noise very
effectively.

The Gaussian filter performs well~(58.0\,\%) because the noise is predominantly
Gaussian in character. It outperforms the moving average despite appearing
superficially similar because the Gaussian kernel is theoretically optimal for
removing Gaussian noise with minimum spectral distortion.

The ML denoisers achieve reductions of only~27--32\,\%, substantially lower than
the Kalman and Gaussian filters. Three factors explain this. First, the training
target is the median-filtered signal rather than a physical thermocouple
measurement, which limits the quality of the learning target. Second, the
training dataset is small~(approximately 1,652 frames from one layer only).
Third, the Kalman filter benefits from physical knowledge encoded in the
$Q$ and $R$ parameters, which the ML models must infer purely from
data~\cite{zhang2023physics}.

\section{Compression Results on NIST Dataset (ATP-3)}
\label{sec:nist-compress}

\subsection{Quantitative Results}

Table~\ref{tab:compress-results} reports the compression ratio and reconstruction
RMSE for all six compression methods applied to the calibrated Layer~01 signal.

\begin{table}[H]
\centering
\caption{Compression results on NIST Layer~01 calibrated signal. Higher ratio
         and lower RMSE indicate better performance. ATP-3 target:
         $CR > 4\times$.}
\label{tab:compress-results}
\begin{tabular}{llcc}
\toprule
\textbf{Method} & \textbf{Type} &
\textbf{CR} & \textbf{RMSE (\textdegree C)} \\
\midrule
Raw Baseline           & Baseline  & $1.0\times$  &   0.00 \\
Downsampling 2$\times$ & Classical & $2.0\times$  &  21.97 \\
SVD rank=10            & Classical & $2.1\times$  & 131.60 \\
Wavelet~(Haar, 10\,\%) & Classical & $5.0\times$  &  36.55 \\
PCA~($n=3$)            & ML        & $10.7\times$ &  58.10 \\
Autoencoder~(bn=3)     & ML        & $10.7\times$ & 170.80 \\
\bottomrule
\end{tabular}
\end{table}

Wavelet compression achieves the best accuracy--ratio trade-off:
$5.0\times$ ratio with RMSE~36.55\,\textdegree C, satisfying the ATP-3 target
of~$CR > 4\times$. PCA achieves the highest ratio~($10.7\times$) at greater
reconstruction error~(58.10\,\textdegree C). The Autoencoder reaches the same
ratio as PCA but with a very high RMSE of~170.80\,\textdegree C. Downsampling
and SVD do not satisfy the ATP-3 ratio target.

Figure~\ref{fig:compress-scatter} shows the ratio versus RMSE trade-off for
all six methods.

\begin{figure}[H]
\centering
\includegraphics[width=\textwidth]{figures/full_comparison_compress}
\caption{Compression ratio versus reconstruction RMSE for all six methods on
         NIST Layer~01. The vertical dashed line marks the ATP-3 target
         of~$4\times$. Wavelet compression offers the best classical trade-off;
         PCA offers the best ML trade-off.}
\label{fig:compress-scatter}
\end{figure}

\subsection{Discussion}

Wavelet compression performs best because the Haar step-function basis is a
natural match for the piecewise-smooth shape of a pyrometer cooling
curve~\cite{learnable_wavelet}. Most of the signal energy is concentrated in a
small number of large-magnitude coefficients at coarse wavelet scales, and
retaining only the top~10\,\% captures dominant behaviour with relatively
low RMSE.

Downsampling achieves the lowest RMSE~(21.97\,\textdegree C) but fails the
ratio target at only~$2\times$. Truncated SVD achieves poor RMSE
(131.6\,\textdegree C) at only~$2.1\times$, because sinusoidal SVD basis
functions require many components to approximate the sharp quench
transition~\cite{golub2013matrix}.

PCA achieves the highest ratio~($10.7\times$) with RMSE~58.10\,\textdegree C.
The Autoencoder reaches the same ratio but with much worse reconstruction error
(170.80\,\textdegree C), likely because the small training dataset of~2,065
frames is insufficient for the network to learn a compact yet accurate
representation~\cite{autoencoder_compression}.

\section{Results on PODFAM Dataset}
\label{sec:podfam-results}

The PODFAM dataset consists of ten point cloud scans from an Aconity MIDI laser
powder bed fusion machine with a real two-colour pyrometer. The data was provided
as part of the PODFAM project. Unlike the NIST dataset, PODFAM contains no
thermocouple reference channel. ML methods therefore cannot be retrained on this
data, and only classical denoising and compression methods are applied.

\subsection{PODFAM Denoising Results}

Table~\ref{tab:podfam-denoise} shows the classical denoising results across all
ten PODFAM files. The raw noise level of approximately~5.4\,\textdegree C is
much lower in absolute terms than on NIST~(500\,\textdegree C), but represents
a similar relative proportion of the signal amplitude.

\begin{table}[H]
\centering
\caption{Classical denoising noise standard deviation~(\textdegree C) across
         ten PODFAM files. Raw noise $\approx 5.4$\,\textdegree C.
         ML methods not applied~(no thermocouple reference available).}
\label{tab:podfam-denoise}
\begin{tabular}{lcccccc}
\toprule
\textbf{File} & \textbf{Raw} & \textbf{MovAvg} & \textbf{Median}
& \textbf{SavGol} & \textbf{Gaussian} & \textbf{Kalman} \\
\midrule
000.280 & 5.42 & 1.73 & 1.93 & 2.25 & 1.29 & 0.63 \\
002.030 & 5.41 & 1.72 & 1.92 & 2.24 & 1.28 & 0.62 \\
003.080 & 5.40 & 1.72 & 1.93 & 2.25 & 1.29 & 0.63 \\
004.060 & 5.41 & 1.73 & 1.93 & 2.25 & 1.29 & 0.63 \\
005.110 & 5.39 & 1.71 & 1.91 & 2.22 & 1.28 & 0.63 \\
006.090 & 5.40 & 1.72 & 1.92 & 2.24 & 1.29 & 0.63 \\
007.000 & 5.39 & 1.72 & 1.92 & 2.23 & 1.29 & 0.63 \\
008.820 & 5.41 & 1.73 & 1.93 & 2.24 & 1.30 & 0.63 \\
009.030 & 5.42 & 1.73 & 1.93 & 2.24 & 1.30 & 0.63 \\
010.010 & 5.40 & 1.72 & 1.92 & 2.24 & 1.29 & 0.63 \\
\midrule
\textbf{Mean}      & 5.41 & 1.72 & 1.92 & 2.24 & 1.29 & 0.63 \\
\textbf{Reduction} & ---  & 68.2\,\% & 64.5\,\% & 58.6\,\%
                          & 76.2\,\% & 88.4\,\% \\
\bottomrule
\end{tabular}
\end{table}

The results are strikingly consistent across all ten files. The Kalman filter
achieves~88.4\,\% noise reduction in every file, reducing the raw noise
from~5.4\,\textdegree C to~0.63\,\textdegree C. The Gaussian filter
achieves~76.2\,\% and the moving average~68.2\,\%. This consistency without any
retuning between files confirms that the classical methods are robust enough for
deployment in an industrial setting where signal conditions vary between
production runs.

\subsection{PODFAM Melt-Pool Temperature Statistics}

Table~\ref{tab:podfam-temps} reports the melt-pool temperature statistics across
all ten PODFAM files after two-colour formula conversion. Temperatures range from
1519 to 1868\,\textdegree C, which is physically plausible for laser powder
bed fusion of metal.

\begin{table}[H]
\centering
\caption{Melt-pool temperature statistics from two-colour pyrometer formula
         across ten PODFAM files. Temperatures in \textdegree C.}
\label{tab:podfam-temps}
\begin{tabular}{lcccc}
\toprule
\textbf{File} & \textbf{Points} & \textbf{Median T}
& \textbf{Max T} & \textbf{Min T} \\
\midrule
000.280 & 2,760,696 & 1867.9 & 2016.2 & 1446.9 \\
002.030 & 2,782,728 & 1594.2 & 1654.0 & 1525.6 \\
003.080 & 2,723,832 & 1661.4 & 1796.0 & 1366.2 \\
004.060 & 2,733,552 & 1628.6 & 1699.4 & 1537.0 \\
005.110 & 2,734,080 & 1519.2 & 1733.4 & 1457.7 \\
006.090 & 2,723,328 & 1725.5 & 1839.5 & 1674.9 \\
007.000 & 2,733,072 & 1554.7 & 1647.3 & 1469.8 \\
008.820 & 2,735,112 & 1652.4 & 1762.5 & 1565.0 \\
009.030 & 2,735,112 & 1584.8 & 1752.7 & 1501.7 \\
010.010 & 2,733,576 & 1650.2 & 1715.8 & 1588.9 \\
\midrule
\textbf{Mean} & 2,741,527 & 1643.9 & 1764.7 & 1513.4 \\
\bottomrule
\end{tabular}
\end{table}

\subsection{PODFAM Compression Results}

Table~\ref{tab:podfam-compress} shows classical compression results across the
ten PODFAM files.

\begin{table}[H]
\centering
\caption{Classical compression results across ten PODFAM files.
         RMSE values in \textdegree C.}
\label{tab:podfam-compress}
\begin{tabular}{lcccccc}
\toprule
& \multicolumn{2}{c}{\textbf{Delta Encoding}}
& \multicolumn{2}{c}{\textbf{Downsampling}}
& \multicolumn{2}{c}{\textbf{Wavelet 10\,\%}} \\
\cmidrule(lr){2-3}\cmidrule(lr){4-5}\cmidrule(lr){6-7}
\textbf{File} & \textbf{CR} & \textbf{RMSE}
              & \textbf{CR} & \textbf{RMSE}
              & \textbf{CR} & \textbf{RMSE} \\
\midrule
000.280 & 2.0 & 0.00 & 2.0 & 0.81 & 3.1 & 1.62 \\
002.030 & 2.0 & 0.00 & 2.0 & 0.81 & 3.1 & 1.69 \\
003.080 & 2.0 & 0.00 & 2.0 & 0.80 & 3.1 & 1.77 \\
004.060 & 2.0 & 0.00 & 2.0 & 0.80 & 3.0 & 1.63 \\
005.110 & 2.0 & 0.00 & 2.0 & 0.80 & 3.1 & 1.67 \\
006.090 & 2.0 & 0.00 & 2.0 & 0.80 & 3.0 & 1.59 \\
007.000 & 2.0 & 0.00 & 2.0 & 0.80 & 3.1 & 1.83 \\
008.820 & 2.0 & 0.00 & 2.0 & 0.81 & 3.0 & 1.59 \\
009.030 & 2.0 & 0.00 & 2.0 & 0.80 & 3.1 & 1.71 \\
010.010 & 2.0 & 0.00 & 2.0 & 0.80 & 3.0 & 1.62 \\
\midrule
\textbf{Mean} & 2.0 & 0.00 & 2.0 & 0.80 & 3.1 & 1.67 \\
\bottomrule
\end{tabular}
\end{table}

Delta encoding achieves near-lossless compression~(RMSE~$\approx 0.0$) at
$2\times$. Downsampling achieves RMSE~0.80\,\textdegree C at~$2\times$.
Wavelet compression achieves~$3.1\times$ with RMSE below~2\,\textdegree C
across all ten files. The wavelet ratio on PODFAM~($3.1\times$) is lower than
on NIST~($5.0\times$) because the PODFAM signal has higher spatial variability
across scan points, requiring more coefficients to represent it accurately.

Figure~\ref{fig:podfam-single} shows the pipeline output on a single PODFAM
file. Figure~\ref{fig:podfam-all} shows the comparison across all ten files.

\begin{figure}[H]
\centering
\includegraphics[width=\textwidth]{figures/podfam_result}
\caption{PODFAM classical pipeline output on file~000.280. Top panels: raw
         versus classical denoised signals. Bottom panels: compression results.
         The Kalman filter~(top right) gives the cleanest output. Results are
         from the real two-colour pyrometer data from the PODFAM project.}
\label{fig:podfam-single}
\end{figure}

\begin{figure}[H]
\centering
\includegraphics[width=\textwidth]{figures/podfam_all_results}
\caption{PODFAM classical pipeline results across all ten files. Each column
         corresponds to one file. The top row shows noise standard deviation
         before and after each denoising method. The bottom row shows
         compression ratio and RMSE. Results are highly consistent across
         all ten files, confirming robustness on real industrial pyrometer data.}
\label{fig:podfam-all}
\end{figure}

\section{Discussion}
\label{sec:discussion}

\subsection{Classical Methods versus ML Methods}

Table~\ref{tab:classical-vs-ml} summarises the head-to-head comparison between
the best classical and best ML method for each stage.

\begin{table}[H]
\centering
\caption{Best classical versus best ML method for each pipeline stage.}
\label{tab:classical-vs-ml}
\begin{tabular}{lllll}
\toprule
\textbf{Stage} & \textbf{Best Classical} & \textbf{Value}
               & \textbf{Best ML} & \textbf{Value} \\
\midrule
Denoising (NIST)   & Kalman filter & 89.5\,\% & LSTM & 31.8\,\% \\
Denoising (PODFAM) & Kalman filter & 88.4\,\% & N/A  & --- \\
Compression (NIST) & Wavelet 10\,\% & $5.0\times$/36.6\textdegree C
                   & PCA $n=3$ & $10.7\times$/58.1\textdegree C \\
\bottomrule
\end{tabular}
\end{table}

For denoising, the classical Kalman filter clearly outperforms all ML methods on
both datasets. The gap is large~(89.5\,\% vs~31.8\,\% on NIST) and is primarily
explained by the quality of the ML training target and the limited training data
size. For compression, neither classical nor ML provides an unambiguous winner:
wavelet compression gives lower RMSE at a lower ratio, while PCA gives a higher
ratio at higher RMSE. The right choice depends on the downstream application's
priorities.

\subsection{Combined Pipeline Performance}

Table~\ref{tab:pipeline-summary} summarises the recommended method for each
pipeline stage based on the experimental results.

\begin{table}[H]
\centering
\caption{Recommended pipeline configuration and key performance metrics.}
\label{tab:pipeline-summary}
\begin{tabular}{llll}
\toprule
\textbf{Stage} & \textbf{Recommended Method}
& \textbf{Metric} & \textbf{Value} \\
\midrule
ATP-1~(real-time)   & Kalman filter   & Noise reduction & 89.5\,\% \\
ATP-1~(offline)     & CNN denoiser    & Noise reduction & 30.0\,\% \\
ATP-2~(calibration) & Random Forest~\cite{ajay2026}
                                      & RMSE vs TC & 108.9\,\textdegree C \\
ATP-3~(accuracy)    & Wavelet 10\,\%  & CR / RMSE
                    & $5.0\times$ / 36.6\,\textdegree C \\
ATP-3~(storage)     & PCA $n=3$       & CR / RMSE
                    & $10.7\times$ / 58.1\,\textdegree C \\
\bottomrule
\end{tabular}
\end{table}

\subsection{Answers to Research Questions}

\textbf{Research Question~1.} The Kalman filter provides the highest noise
reduction of all ten methods tested: 89.5\,\% on NIST and 88.4\,\% on PODFAM.
The CNN achieves the best ML result~(30.0\,\%). The median filter is recommended
for industrial deployment because it removes impulse spikes reliably on both
datasets without requiring any training data while preserving peak temperatures
accurately. ATP-1 is satisfied --- all methods achieve a demonstrably lower noise
standard deviation than the raw signal baseline.

\textbf{Research Question~2.} Wavelet compression with 10\,\% Haar coefficient
retention achieves a~$5.0\times$ ratio with RMSE~36.55\,\textdegree C,
satisfying the ATP-3 target of~$CR > 4\times$. PCA achieves the highest
ratio~($10.7\times$) at greater reconstruction error~(58.10\,\textdegree C).
Downsampling and SVD do not satisfy the ratio target. ATP-3 is satisfied by
the wavelet method, and for applications tolerating higher reconstruction error,
by PCA compression.

\let\cleardoublepage\clearpage
